% !TEX program = xelatex
% ============================================================
%  متتبع الميزانية — Budget Tracker Template
%  نموذج لتتبع ميزانية المشروع ومقارنة المخطط بالفعلي
%  Compile with: xelatex main.tex (run twice)
% ============================================================
%  Features:
%    - Budget table: Category | Budgeted | Actual | Variance | % used
%    - Subtotals and total row
%    - Monthly burn table
% ============================================================

\documentclass[11pt,a4paper]{article}

\usepackage[a4paper,margin=2.5cm,top=2.5cm,bottom=2.5cm,headheight=16pt]{geometry}
\usepackage{graphicx}
\usepackage{array}
\usepackage{booktabs}
\usepackage{tabularx}
\usepackage{multirow}
\usepackage{enumitem}
\usepackage{float}
\usepackage{titlesec}
\usepackage{fancyhdr}
\usepackage{setspace}
\usepackage{xcolor}
\usepackage{amsmath}

\usepackage{bidi}
\usepackage[no-math]{fontspec}
\usepackage[quiet]{polyglossia}

\setmainfont[Scale=1]{Amiri-Regular.ttf}
\newfontfamily\arabicfont[Script=Arabic,Scale=0.95]{Amiri-Regular.ttf}
\newfontfamily\englishfont[Scale=0.95]{TeX Gyre Termes}

\setdefaultlanguage[calendar=gregorian,hijricorrection=1]{arabic}
\setotherlanguage[variant=british]{english}

\usepackage[breaklinks,hidelinks]{hyperref}

\addto\captionsarabic{%
  \renewcommand{\contentsname}{الفهرس}%
  \renewcommand{\figurename}{الشكل}%
  \renewcommand{\tablename}{الجدول}%
}

\titleformat{\section}{\Large\bfseries}{\thesection}{0.7em}{}
\titlespacing*{\section}{0pt}{1.4ex plus 0.4ex}{0.9ex plus 0.3ex}

\pagestyle{fancy}
\fancyhf{}
\fancyhead[R]{\small\itshape متتبع الميزانية}
\fancyhead[L]{\small اسم المشروع}
\fancyfoot[C]{\small\thepage}
\renewcommand{\headrulewidth}{0.3pt}

\setlist[itemize]{topsep=0.3em, itemsep=0.15em, parsep=0pt}
\setlist[enumerate]{topsep=0.3em, itemsep=0.15em, parsep=0pt}

\newcolumntype{L}[1]{>{\raggedright\arraybackslash}p{#1}}
\newcolumntype{C}[1]{>{\centering\arraybackslash}p{#1}}
\newcolumntype{R}[1]{>{\raggedleft\arraybackslash}p{#1}}

\setstretch{1.3}

\begin{document}

% ============================================================
% TITLE BLOCK
% ============================================================
\begin{center}
{\LARGE\bfseries متتبع الميزانية}\\[0.3em]
{\large Budget Tracker}\\[1em]
\end{center}

% ============================================================
% PROJECT INFO
% ============================================================
% TODO: املأ بيانات المشروع.
\begin{table}[H]
\centering
\renewcommand{\arraystretch}{1.4}
\begin{tabularx}{\textwidth}{L{3.5cm} X L{3.5cm} X}
\toprule
\textbf{اسم المشروع:} & --- & \textbf{مدير المشروع:} & --- \\
\textbf{اسم الشركة:} & --- & \textbf{تاريخ التحديث:} & --- \\
\textbf{العملة:} & \LR{USD} & \textbf{إجمالي الميزانية:} & --- \\
\bottomrule
\end{tabularx}
\end{table}

\vspace{1em}

% ============================================================
% 1. BUDGET BY CATEGORY
% ============================================================
\section{الميزانية حسب الفئة}
\label{sec:budget}

% TODO: عدّل الفئات والمبالغ حسب ميزانية مشروعك.
% Variance = Budgeted - Actual (positive = under budget).
\begin{table}[H]
\centering
\caption{الميزانية حسب الفئة}
\renewcommand{\arraystretch}{1.4}
\begin{tabularx}{\textwidth}{L{4cm} R{3cm} R{3cm} R{3cm} C{2cm}}
\toprule
\textbf{الفئة} & \textbf{المخطط} & \textbf{الفعلي} & \textbf{الانحراف} & \textbf{\% مستخدم} \\
\midrule
\multicolumn{5}{l}{\textbf{الرواتب والموارد البشرية}} \\
\midrule
الرواتب والأجور & 50{,}000 & 48{,}000 & 2{,}000 & 96\% \\
الاستشارات الخارجية & 15{,}000 & 12{,}000 & 3{,}000 & 80\% \\
\midrule
\textbf{المجموع الفرعي} & \textbf{65{,}000} & \textbf{60{,}000} & \textbf{5{,}000} & \textbf{92\%} \\
\midrule
\multicolumn{5}{l}{\textbf{المعدات والمواد}} \\
\midrule
المعدات والأجهزة & 20{,}000 & 22{,}000 & -2{,}000 & 110\% \\
المواد الاستهلاكية & 5{,}000 & 4{,}500 & 500 & 90\% \\
\midrule
\textbf{المجموع الفرعي} & \textbf{25{,}000} & \textbf{26{,}500} & \textbf{-1{,}500} & \textbf{106\%} \\
\midrule
\multicolumn{5}{l}{\textbf{التشغيل والخدمات}} \\
\midrule
السفر والمؤتمرات & 8{,}000 & 6{,}500 & 1{,}500 & 81\% \\
النشر والتوثيق & 4{,}000 & 3{,}000 & 1{,}000 & 75\% \\
خدمات الاتصالات & 3{,}000 & 3{,}000 & 0 & 100\% \\
\midrule
\textbf{المجموع الفرعي} & \textbf{15{,}000} & \textbf{12{,}500} & \textbf{2{,}500} & \textbf{83\%} \\
\midrule
\multicolumn{5}{l}{\textbf{احتياطي}} \\
\midrule
احتياطي الطوارئ (10\%) & 10{,}500 & 2{,}000 & 8{,}500 & 19\% \\
\midrule
\textbf{الإجمالي} & \textbf{115{,}500} & \textbf{101{,}000} & \textbf{14{,}500} & \textbf{88\%} \\
\bottomrule
\end{tabularx}
\end{table}

\vspace{0.5em}

\noindent\textbf{ملاحظة:} الانحراف الموجب يعني أن الإنفاق أقل من المخطط (تحت الميزانية)، والانحراف السالب يعني تجاوز الميزانية.

% ============================================================
% 2. MONTHLY BURN
% ============================================================
\section{الاستهلاك الشهري للميزانية}
\label{sec:burn}

% TODO: عدّل بيانات الاستهلاك الشهري حسب مشروعك.
\begin{table}[H]
\centering
\caption{الاستهلاك الشهري للميزانية}
\renewcommand{\arraystretch}{1.4}
\begin{tabularx}{\textwidth}{C{2cm} R{3cm} R{3cm} R{3cm} C{2.5cm}}
\toprule
\textbf{الشهر} & \textbf{المخطط الشهري} & \textbf{الفعلي الشهري} & \textbf{المتراكم} & \textbf{\% من الإجمالي} \\
\midrule
الشهر 1 & 10{,}000 & 9{,}500 & 9{,}500 & 8\% \\
الشهر 2 & 12{,}000 & 11{,}000 & 20{,}500 & 18\% \\
الشهر 3 & 15{,}000 & 14{,}000 & 34{,}500 & 30\% \\
الشهر 4 & 15{,}000 & 16{,}500 & 51{,}000 & 44\% \\
الشهر 5 & 18{,}000 & 17{,}000 & 68{,}000 & 59\% \\
الشهر 6 & 20{,}000 & 18{,}500 & 86{,}500 & 75\% \\
الشهر 7 & 15{,}500 & 14{,}500 & 101{,}000 & 88\% \\
\midrule
\textbf{الإجمالي} & \textbf{105{,}500} & \textbf{101{,}000} & --- & \textbf{88\%} \\
\bottomrule
\end{tabularx}
\end{table}

% ============================================================
% 3. NOTES
% ============================================================
\section{ملاحظات}
\label{sec:notes}

% TODO: أضف ملاحظات حول الميزانية.
\begin{itemize}
    \item فئة المعدات تجاوزت الميزانية بنسبة 10\% بسبب زيادة أسعار الأجهزة.
    \item الاستهلاك الكلي أقل من المخطط بنسبة 12\%، مما يوفر هامشاً مالياً إيجابياً.
    \item احتياطي الطوارئ لم يُستخدم بالكامل بعد، ويمكن إعادة تخصيصه إذا لزم الأمر.
\end{itemize}

\end{document}
